之间通常认为是一个大的样本. 随着数据特征的变大,为了得到不同的检验统计量的真实分布的良好的渐进逼近,我们需要更大的样本容量. 我们不打算论证地简单认为 $n-p$ 必须是大的,并且意识到真实情况会更加复杂. 对 $p=2$ 且样本容量为 50 的应用和 $p=52$ 且样本容量为 100 的应用是很不一样,尽管两种情况 $n-p$ 都等于 48.

对这些大样本推断方法进行正态理论所要求的各种检验,这是一种很好的统计实践. 虽然当 $n$ 很大时,对正态性的微小甚至中等程度的偏离一般不会引起任何问题,但过大的偏差就会产生问题,特别是真实误差率会远远地偏离其名义误差率水平 $\alpha$. 如果在 Q-Q 散点图和其他考察方法中发现存在离群值和其他过大的偏差(例如,参见[2]),则应采取适当的校正措施,包括对数据进行变换. 对称多元分布均值向量的检验方法相对来说,对偏离正态性不很敏感,这些方法在文献[11]中有所讨论. 在某些情形下,结论 5.4 和结论 5.5 仅对非常大的样本才有效.

下面的例子用来解释如何对单一均值向量构造大样本联合置信区间的问题.

\noindent\textbf{例 5.7}\quad (构造大样本联合置信区间)

一名音乐教育工作者为制定芬兰的音乐教育标准,测试了数千名芬兰学生的音乐才能,部分统计数据见\hyperref[tab:5.5]{表 5.5}. 这些统计数据来自自由 $n=96$ 名芬兰 12 年级学生组成的样本:

\begin{mvtable}
\phantomsection\label{tab:5.5}\mvtabcap{表 5.5\quad 96 名 12 年级芬兰学生音乐才能测试数据的均值与标准差}
\begin{tabular}{lcc}
\toprule
\multirow{2}{*}{变量} & \multicolumn{2}{c}{原始成绩}\\
\cmidrule(lr){2-3}
 & 均值($\bar{x}_i$) & 标准差($\sqrt{s_{ii}}$)\\
\midrule
$X_1$=旋律 & 28.1 & 5.76\\
$X_2$=和声 & 26.6 & 5.85\\
$X_3$=速度 & 35.4 & 3.82\\
$X_4$=拍子 & 34.2 & 5.12\\
$X_5$=短句 & 23.6 & 3.76\\
$X_6$=协调 & 22.0 & 3.93\\
$X_7$=风格 & 22.7 & 4.03\\
\bottomrule
\end{tabular}
\end{mvtable}

资料来源:数据经 V.\ Sell 许可使用.

我们来构造均值分量 $\mu_i,i=1,2,\cdots,7$ 的 $90\%$ 联合置信区间.

由结论 5.5,$90\%$ 的联合置信限为
\[
\bar{x}_i\pm\sqrt{\chi_7^{2}(0.10)}\ \sqrt{\frac{s_{ii}}{n}},\qquad i=1,2,\cdots,7,\ \text{其中}\ \chi_7^{2}(0.10)=12.02
\]
因此在大约 $90\%$ 的置信水平下有:
\begin{align*}
28.1\pm\sqrt{12.02}\times\frac{5.76}{\sqrt{96}}&\quad\text{包含}\ \mu_1,\ \text{即}\ 26.06\leqslant\mu_1\leqslant30.14\\
26.6\pm\sqrt{12.02}\times\frac{5.85}{\sqrt{96}}&\quad\text{包含}\ \mu_2,\ \text{即}\ 24.53\leqslant\mu_2\leqslant28.67\\
35.4\pm\sqrt{12.02}\times\frac{3.82}{\sqrt{96}}&\quad\text{包含}\ \mu_3,\ \text{即}\ 34.05\leqslant\mu_3\leqslant36.75\\
34.2\pm\sqrt{12.02}\times\frac{5.12}{\sqrt{96}}&\quad\text{包含}\ \mu_4,\ \text{即}\ 32.39\leqslant\mu_4\leqslant36.01
\end{align*}
